% =====================================================================
%  Thème 3 — Colinéarité et parallélisme — Niveau FACILE — Corrigé
% =====================================================================
\documentclass[11pt,a4paper]{article}
\usepackage[utf8]{inputenc}
\usepackage[T1]{fontenc}
\usepackage{lmodern}
\usepackage[french]{babel}
\usepackage{amsmath,amssymb,mathtools}
\usepackage{xcolor}
\usepackage{enumitem}
\usepackage{fancyhdr}
\usepackage[a4paper,margin=2.1cm,headheight=26pt,headsep=10pt]{geometry}

\definecolor{primary}{RGB}{222,70,80}
\definecolor{primarydark}{RGB}{150,30,42}
\definecolor{excol}{RGB}{0,135,90}

\pagestyle{fancy}\fancyhf{}
\fancyhead[L]{\small\textcolor{primarydark}{\textbf{Fiche 2} — Calcul vectoriel}}
\fancyhead[R]{\small\textcolor{primarydark}{\textsc{Thème 3 — Facile — Corrigé}}}
\fancyfoot[C]{\small\thepage}
\renewcommand{\headrulewidth}{0.8pt}

\newcommand{\vc}[1]{\overrightarrow{#1}}
\providecommand{\norm}[1]{\left\lVert #1\right\rVert}
\newcommand{\cor}[1]{\medskip\noindent{\color{primarydark}\bfseries\large Corrigé #1.}\par\smallskip}
\newcommand{\rep}{\textcolor{excol}{$\blacktriangleright$}\ }

\begin{document}
\begin{center}
{\LARGE\bfseries\color{primarydark}Colinéarité et parallélisme — Corrigé Facile}\\[2pt]
{\color{primary}\rule{0.5\linewidth}{1.1pt}}
\end{center}
\vspace{2mm}

\cor{1}
\begin{enumerate}[label=\textbf{\alph*)},leftmargin=2em,itemsep=2pt]
  \item $\vec v=(6;9)=3\,(2;3)=3\vec u$ : colinéaires, $k=3$.
  \item $\vec v=(-3;6)=-3\,(1;-2)=-3\vec u$ : colinéaires, $k=-3$.
  \item $\vec v=(4;5)$ : il faudrait $4=2k$ (donc $k=2$) \emph{et} $5=3k$ (donc $k=\frac53$).
        Impossible : \textbf{non} colinéaires.
\end{enumerate}
\rep Deux vecteurs sont colinéaires si \emph{un seul et même} $k$ marche pour les deux composantes.

\cor{2}
\begin{enumerate}[label=\textbf{\alph*)},leftmargin=2em,itemsep=2pt]
  \item $3\times8-6\times4=24-24=0$ : \textbf{colinéaires}.
  \item $1\times5-3\times2=5-6=-1\neq0$ : \textbf{non} colinéaires.
  \item $(-2)\times(-10)-4\times5=20-20=0$ : \textbf{colinéaires}.
\end{enumerate}
\rep Le déterminant se calcule « en croix » : $x\cdot y'$ moins $x'\cdot y$.

\cor{3}
\begin{enumerate}[label=\textbf{\alph*)},leftmargin=2em,itemsep=2pt]
  \item $\vc{AB}\,(2;1)$, $\vc{AC}\,(6;3)$ ; déterminant $=2\times3-6\times1=0$ : $A,B,C$ \textbf{alignés}.
  \item $\vc{AB}\,(2;3)$, $\vc{AC}\,(4;5)$ ; déterminant $=2\times5-4\times3=10-12=-2\neq0$ :
        \textbf{non} alignés (vrai triangle).
\end{enumerate}

\cor{4}
\begin{enumerate}[label=\textbf{\alph*)},leftmargin=2em,itemsep=2pt]
  \item $a=2,\,b=3$ : $\vec u\,(-3;2)$.
  \item $a=1,\,b=-4$ : $\vec u\,(4;1)$.
  \item $y=2x+1 \Rightarrow 2x-y+1=0$, donc $a=2,\,b=-1$ : $\vec u\,(1;2)$.
\end{enumerate}
\rep On lit $a$ et $b$ dans $ax+by+c=0$, puis $\vec u\,(-b;a)$.

\cor{5}
\begin{enumerate}[label=\textbf{\alph*)},leftmargin=2em,itemsep=2pt]
  \item $m=\dfrac{y_B-y_A}{x_B-x_A}=\dfrac{8-2}{4-1}=\dfrac63=2$.
  \item $\vec u\,(3;6)$ : $m=\dfrac{6}{3}=2$.
  \item $3x+2y-6=0$ : $m=-\dfrac{a}{b}=-\dfrac{3}{2}$.
\end{enumerate}
\rep La pente est « variation de $y$ sur variation de $x$ », soit $\frac{y_B-y_A}{x_B-x_A}$ ou, à partir
d'un vecteur directeur $(u_1;u_2)$, le rapport $\frac{u_2}{u_1}$.

\end{document}
